\section{Conclusion and Open Problems}\label{sec:conclusion}

In the model, seven pawns on four files of the ordinary board make a $\nim{2}$ and no six do; nine make a
$\nim{3}$. On $2{,}049$ decisive $8\times8$ positions that we could solve (not a random sample) the model names the real winner in $92.5\%$; on six ranks a
square-rule passer test removes its pawn errors with frozen kings \lab{C1}. Section~S10 lists two more problems.

\begin{enumerate}
\item \emph{Higher nimbers.} Can a ring $\nim{3}$ on $\cyl{n}$ have only $n$ pawns (Fig.~\ref{fig:least})? Are
  the wall families $\nim{2}$ as sums for all $n$? Is there a $\nim{4}$ at any height?
\item \emph{Real chess.} With queen promotions and frozen kings the oracle value of the $1/8$ position is exactly
  $1/8$; does it survive all four promotions? The 8-man end positions deciding A and B agree with the model;
  those with 9 to 11 men are unsolved.
\item \emph{When is a tr\'ebuchet zero?} The certificate theory settles the promotion game, not chess. In the promotion
  game even a locked pawn pair can break a tr\'ebuchet's zero \lab{C2} (Section~S2).
\item \emph{Mis\`ere Dawson's chess.} Dawson asked about mis\`ere play \cite{Dawson1935}: is the mis\`ere
  quotient \cite{PlambeckSiegel2008} of Dawson's chess infinite? By the ring rule (Theorem~\ref{thm:ring}), an
  answer for strips also answers the cylinder.
\end{enumerate}
